% page 576 of the facsimile (image p-575 of the scan)
% block 160.0 x 254.9 mm ; left margin 8.0 mm
\pgc{-12.700mm}{0.000mm}{
\l{}
\l{}
\l{}
\l{}
\l{\cel{3}568}
\l{}
\l{\cel{3}felo tensita, sur un ek qui onu frapas per vergeto por igar}
\l{\cel{3}lu resonar. - EFIRS.}
\l{}
\l{\sou{tamburino}. Tamburo plu basa kam la tamburo ordinara, quan onu}
\l{\cel{3}igas resonar per nur un vergeto. -}
\l{}
\l{\sou{tamen}. (\sou{Konjunciono})\cel{2}Opoze ad ico ; malgre lo dicita ; malgre}
\l{\cel{3}ta obstaklo.- L.}
\l{}
\l{\sou{tampar}. (\sou{trans.}) Igar (pavo, paliso, e c.) irar ad, o til,}
\l{\cel{3}la fundo per frapar (lu) per grava bloko ek ligno. - E.}
\l{}
\l{\sou{tampono}. I. Mikra bloko, garnisita ye linjo, felto, e c., quan}
\l{\cel{3}onu imbibas per substanco liquida quan onu volas extensar,}
\l{\cel{3}aplikar sur surfaco. - II. (\sou{skermarto}). Rondo de felo quan}
\l{\cel{3}onu muntas tampone ye la extremajo di flereto por igar ica}
\l{\cel{3}nevundiva. - III. (\sou{tekn.)}\cel{2}Stifto ligna, ye qua onu garni-}
\l{\cel{3}sas truo borita en muro, por fixigar en lu skrubo, klovo.}
\l{\cel{3}- DEFS.}
\l{}
\l{\sou{tamtamo}. Perkut-instrumento (che la populi di la esto-regiono}
\l{\cel{3}di Azia), disko metala, kun rebordo, qua sonegas kande onu}
\l{\cel{3}frapas lu. DEFIR.}
\l{}
\l{\sou{tano}. Kortico di querko, di sumako, di kastaniero, pulverigi-}
\l{\cel{3}ta, quan onu uzas por la preparo di la ledri. - EF.}
\l{}
\l{\sou{tandem}. Indikas ke lo eventis pos lasir (ulu)expektar ta even-}
\l{\cel{3}to dum multa tempo. - L.}
\l{}
\l{\sou{tandemo}. I. Veturo kun du kavali jungita an singla latero di}
\l{\cel{3}timono. - II. Biciklo por du personi qui kavalkas lu, la}
\l{\cel{3}duesma dop la unesma. - DEFI.}
\l{}
\l{tangar. (\sou{netrans}.) (\sou{Navig}.) (\sou{aludante navo an aquo}). Ocilar al-}
\l{\cel{3}tarne de la pruo a la pupo, de la pupo a la pruo. - F.}
\l{}
\l{\sou{tangenta}. (\sou{geometrio}). Qua tushas lineo, surfaco, ye nur un}
\l{\cel{3}punto. - DEFIRS.}
\l{}
\l{\sou{tanino}. (\sou{kemio)}. Substanco astriktiva, qua existas en la kor-}
\l{\cel{3}tico di querko, sumako, kastaniero, ed igas olu apta pre-}
\l{\cel{3}parar la ledri. DEFIRS.}
\l{}
\l{\sou{tanko}. Granda recipiento metala, klozita e klozebla hermetike,}
\l{\cel{3}por la liquidi. - DEFS.}
\l{}
\l{\sou{tanta}. Indikas quanto, grado, nombro, nedefinita o tre granda.}
\l{\cel{3}FIS.}
\l{}
\l{\sou{tantalo.}\cel{1}(\sou{kemio)}. Metalo duktila, tre elastika, fasonebla por}
\l{\cel{3}la faco di skribo-plumi, siringi por injekti,subderma, e c.}
\l{\cel{3}-\cel{1}\sou{Simbolo kemiala}\cel{1}: Ta. - DEFIRS}
}
